\documentclass[aspectratio=169,10pt]{beamer}
\usetheme[numbering=fraction,progressbar=frametitle]{metropolis}
\usepackage{fontspec}
\setsansfont{Latin Modern Sans}
\usepackage{amsmath,amssymb,mathtools,graphicx,booktabs}
\usepackage{polyglossia}
\setdefaultlanguage{french}
\setbeamertemplate{navigation symbols}{}
\setbeamerfont{frametitle}{size=\large}
\setbeamerfont{normal text}{size=\normalsize}
\setlength{\parskip}{.45em}
\hypersetup{pdfauthor={},pdfsubject={Lecture slides},pdfkeywords={},pdftitle={Nombres complexes : point de vue géométrique}}
\providecommand{\N}{\mathbb{N}}
\providecommand{\Z}{\mathbb{Z}}
\providecommand{\D}{\mathbb{D}}
\providecommand{\Q}{\mathbb{Q}}
\providecommand{\R}{\mathbb{R}}
\providecommand{\C}{\mathbb{C}}
\providecommand{\vect}[1]{\overrightarrow{#1}}
\providecommand{\vi}{\vec{\imath}}
\providecommand{\vj}{\vec{\jmath}}
\providecommand{\vk}{\vec{k}}
\providecommand{\norm}[1]{\left\lVert #1\right\rVert}
\providecommand{\vc}[2]{\begin{pmatrix}#1\\#2\end{pmatrix}}
\providecommand{\deter}[2]{\det\left(#1,#2\right)}
\providecommand{\abs}[1]{\left\lvert #1\right\rvert}
\providecommand{\dd}{\,\mathrm{d}}
\providecommand{\e}{\mathrm{e}}
\providecommand{\Card}{\operatorname{Card}}
\providecommand{\E}{\operatorname{E}}
\providecommand{\V}{\operatorname{V}}
\providecommand{\degre}{^{\circ}}
\providecommand{\un}[1]{\,\mathrm{#1}}
\providecommand{\sci}[2]{#1\times10^{#2}}
\providecommand{\pgcd}{\operatorname{PGCD}}
\providecommand{\ppcm}{\operatorname{PPCM}}
\title{Nombres complexes : point de vue géométrique}
\subtitle{Cours et exemples}
\date{}
\author{}
\begin{document}
\begin{frame}[plain]
\vfill
{\usebeamerfont{title}\inserttitle\par}
\medskip
{\usebeamerfont{subtitle}\insertsubtitle\par}
\medskip
{\usebeamercolor[fg]{progress bar}\rule{\textwidth}{.5pt}}
\vfill
\end{frame}
\begin{frame}{Objectifs}
Représenter un nombre complexe par un point ou un vecteur ; déterminer un module et un argument ; passer de la forme algébrique à la forme trigonométrique et inversement ; interpréter un module comme une distance.
\end{frame}
\begin{frame}{Activité d'introduction}
Le plan est muni d'un repère orthonormé \((O;\vect u,\vect v)\). À un nombre complexe \(x+iy\) on associe le point de coordonnées \((x;y)\).
\end{frame}
\begin{frame}{Activité d'introduction (2)}
\begin{enumerate}
\setcounter{enumi}{0}\item Placer les points \(A\), \(B\), \(C\) associés à \(a=2+i\), \(b=i\,a\) et \(c=(1+i)\,a\). Donner \(b\) et \(c\) sous forme algébrique.

\item Calculer les distances \(OA\), \(OB\) et \(OC\).

\item Calculer le produit scalaire \(\vect{OA}\cdot\vect{OB}\). Que peut-on dire des droites \((OA)\) et \((OB)\) ?

\item Calculer \(\cos\left(\vect{OA},\vect{OC}\right)\) à l'aide du produit scalaire, puis le signe du déterminant \(\deter{\vect{OA}}{\vect{OC}}\). En déduire une mesure de l'angle orienté \(\left(\vect{OA},\vect{OC}\right)\).

\end{enumerate}
\end{frame}
\begin{frame}{Activité d'introduction (3)}
\begin{enumerate}
\setcounter{enumi}{4}\item Recommencer avec \(a=-1+3i\) (ou un point de votre choix). Formuler une conjecture sur l'effet géométrique d'une multiplication par \(i\), puis par \(1+i\).

\item Quelles grandeurs associées à un nombre complexe faudrait-il définir pour énoncer cette conjecture en général ?

\end{enumerate}
\end{frame}
\begin{frame}{Activité d'introduction (4)}
\emph{Observer n'est pas démontrer : les propriétés conjecturées ici sont établies dans le cours.}
\end{frame}
\begin{frame}{Plan complexe}
Le plan est muni d'un repère orthonormé direct \((O;\vect u,\vect v)\). Au nombre complexe \(z=a+ib\) (\(a\), \(b\) réels) on associe le point \(M(a;b)\), appelé \emph{image} de \(z\), et le vecteur \(\vect w\vc ab\). On dit que \(z\) est l'\emph{affixe} de \(M\) et de \(\vect w\). Le plan muni de cette correspondance est appelé \emph{plan complexe}.
\end{frame}
\begin{frame}{Plan complexe}
Si \(A\) et \(B\) ont pour affixes \(z_A\) et \(z_B\), le vecteur \(\vect{AB}\) a pour affixe \(z_B-z_A\) et le milieu de \([AB]\) a pour affixe \(\dfrac{z_A+z_B}2\). L'affixe de \(\vect w+\vect{w'}\) est la somme des affixes, celle de \(k\vect w\) (\(k\) réel) est \(k\) fois l'affixe de \(\vect w\).

L'image de \(\overline z\) est la symétrique de l'image de \(z\) par rapport à l'axe des abscisses ; l'image de \(-z\) est sa symétrique par rapport à \(O\).
\end{frame}
\begin{frame}{Module}
Le \emph{module} de \(z=a+ib\) est le réel positif \(\abs z=\sqrt{a^2+b^2}\). Si \(M\) est l'image de \(z\), \(\abs z=OM\). Pour un réel, le module coïncide avec la valeur absolue.
\end{frame}
\begin{frame}{Module}
Pour tout complexe \(z\), \(\abs z^2=z\overline z\), et \(\abs{\overline z}=\abs{-z}=\abs z\). Si \(A\) et \(B\) ont pour affixes \(z_A\) et \(z_B\), alors \(AB=\abs{z_B-z_A}\).
\end{frame}
\begin{frame}{Module}
Avec \(z=a+ib\), \(z\overline z=a^2+b^2=\abs z^2\) (chapitre 1). La distance \(AB\) est la norme du vecteur \(\vect{AB}\), d'affixe \(z_B-z_A\).
\end{frame}
\begin{frame}{Module}
Pour tous complexes \(z\), \(z'\) et tout entier naturel \(n\) :
\[\abs{zz'}=\abs z\abs{z'},\qquad\abs{z^n}=\abs z^n,\qquad\text{et, si }z'\neq0,\quad\abs{\dfrac1{z'}}=\dfrac1{\abs{z'}},\quad\abs{\dfrac z{z'}}=\dfrac{\abs z}{\abs{z'}}.\]
\end{frame}
\begin{frame}{Module}
\emph{Produit.} \(\abs{zz'}^2=zz'\,\overline{zz'}=z\overline z\,z'\overline{z'}=\abs z^2\abs{z'}^2\). Les modules sont positifs, donc \(\abs{zz'}=\abs z\abs{z'}\).

\emph{Puissance.} Par récurrence : \(\abs{z^0}=1=\abs z^0\), et si \(\abs{z^n}=\abs z^n\), alors \(\abs{z^{n+1}}=\abs{z^n}\abs z=\abs z^{n+1}\).

\emph{Inverse et quotient.} \(\abs{z'}\times\abs{\dfrac1{z'}}=\abs{z'\times\dfrac1{z'}}=1\), puis \(\abs{\dfrac z{z'}}=\abs z\abs{\dfrac1{z'}}\).
\end{frame}
\begin{frame}{Module}
Pour tous complexes \(z\) et \(z'\), \(\abs{z+z'}\le\abs z+\abs{z'}\). Le problème du chapitre en donne une démonstration et le cas d'égalité.
\end{frame}
\begin{frame}{Module}
On note \(\mathbb U\) l'ensemble des nombres complexes de module \(1\). Leurs images forment le cercle de centre \(O\) et de rayon \(1\).
\end{frame}
\begin{frame}{Module}
Si \(z\) et \(z'\) appartiennent à \(\mathbb U\), alors \(zz'\in\mathbb U\) et \(\dfrac1z\in\mathbb U\) ; de plus \(\dfrac1z=\overline z\).
\end{frame}
\begin{frame}{Module}
\(\abs{zz'}=1\times1=1\) et \(\abs{\dfrac1z}=\dfrac11=1\). Enfin \(z\overline z=\abs z^2=1\), donc \(\overline z=\dfrac1z\).
\end{frame}
\begin{frame}{Module}
Soit \(A\) et \(B\) d'affixes \(a\) et \(b\), et \(M\) d'affixe \(z\). L'égalité \(\abs{z-a}=r\) (\(r>0\)) signifie \(AM=r\) : \(M\) décrit le cercle de centre \(A\) et de rayon \(r\). L'égalité \(\abs{z-a}=\abs{z-b}\) (\(a\neq b\)) signifie \(AM=BM\) : \(M\) décrit la médiatrice de \([AB]\).
\end{frame}
\begin{frame}{Argument et forme trigonométrique}
Soit \(z\) un complexe non nul d'image \(M\). On appelle \emph{argument} de \(z\) toute mesure, en radians, de l'angle orienté \(\left(\vect u,\vect{OM}\right)\). On note \(\arg z\) l'une d'elles ; les autres en diffèrent d'un multiple de \(2\pi\), ce que l'on écrit \(\arg z\equiv\theta\pmod{2\pi}\). Le nombre \(0\) n'a pas d'argument.
\end{frame}
\begin{frame}{Argument et forme trigonométrique}
Tout complexe non nul \(z=a+ib\) s'écrit \(z=r(\cos\theta+i\sin\theta)\) avec \(r=\abs z>0\) et \(\theta\equiv\arg z\pmod{2\pi}\) ; on a alors
\[\cos\theta=\dfrac ar,\qquad\sin\theta=\dfrac br.\]
Réciproquement, si \(z=r(\cos\theta+i\sin\theta)\) avec \(r>0\), alors \(\abs z=r\) et \(\theta\) est un argument de \(z\). Deux formes trigonométriques \(r(\cos\theta+i\sin\theta)\) et \(r'(\cos\theta'+i\sin\theta')\) sont égales si et seulement si \(r=r'\) et \(\theta\equiv\theta'\pmod{2\pi}\).
\end{frame}
\begin{frame}{Argument et forme trigonométrique}
Pour \(z\neq0\) : \(\arg\overline z\equiv-\arg z\) et \(\arg(-z)\equiv\arg z+\pi\pmod{2\pi}\). Un réel strictement positif a pour argument \(0\), un réel strictement négatif \(\pi\) ; \(i\) a pour argument \(\dfrac\pi2\).
\end{frame}
\begin{frame}{Argument et forme trigonométrique}
Calculer \(r=\abs z\), puis \(\dfrac ar\) et \(\dfrac br\), et reconnaître le cosinus et le sinus d'un angle remarquable.
\end{frame}
\begin{frame}{Argument et forme trigonométrique}
\(1+i\sqrt3\) : \(r=2\), \(\cos\theta=\dfrac12\), \(\sin\theta=\dfrac{\sqrt3}2\), donc \(1+i\sqrt3=2\left(\cos\dfrac\pi3+i\sin\dfrac\pi3\right)\).

\(-1+i=\sqrt2\left(\cos\dfrac{3\pi}4+i\sin\dfrac{3\pi}4\right)\) ; \(-3=3(\cos\pi+i\sin\pi)\) ; \(-2i=2\left(\cos\left(-\dfrac\pi2\right)+i\sin\left(-\dfrac\pi2\right)\right)\).
\end{frame}
\begin{frame}{Argument et forme trigonométrique}
\begin{columns}[T,onlytextwidth]
\begin{column}{.57\textwidth}
\(-2\left(\cos\dfrac\pi5+i\sin\dfrac\pi5\right)\) n'est pas une forme trigonométrique, car \(-2<0\). On écrit \(-2\left(\cos\dfrac\pi5+i\sin\dfrac\pi5\right)=2\left(\cos\dfrac{6\pi}5+i\sin\dfrac{6\pi}5\right)\).
\end{column}
\begin{column}{.4\textwidth}
\centering
\includegraphics[width=\linewidth,height=4.5cm,keepaspectratio]{figures/expertes-02-1.pdf}
\end{column}
\end{columns}
\end{frame}
\begin{frame}{Argument et forme trigonométrique}
\(\abs z=OM\), \(\abs{z_B-z_A}=AB\), \(\abs z^2=z\overline z\) ; le module d'un produit est le produit des modules ; \(z=r(\cos\theta+i\sin\theta)\) avec \(r>0\) ; un argument n'est défini qu'à \(2\pi\) près.
\end{frame}
\end{document}
